% ================================================
% 文件: 03_专业测试仪器.tex
% 章节: 第19章 专业测试仪器
% 所属部分: 仪器仪表
% 版本: v1.0
% ================================================

\chapter{专业测试仪器}
\label{chap:specialized_instruments}

\section{频谱分析仪}
\label{sec:spectrum_analyzer}

频谱分析仪是用于分析信号频率成分的仪器。如果说示波器让我们“看见信号随时间如何变化”，那么频谱分析仪则让我们“看见信号在不同频率上各含多少能量”。在射频与通信工程中，发射机的杂散、接收机的本振泄漏、放大器的谐波与互调产物，往往只能在频域中被清晰识别，因此频谱仪是射频实验室与电磁兼容（EMC）测试的核心设备。

\subsection{工作原理}

频谱分析仪将信号分解为不同频率的分量，并显示其幅度谱。现代频谱仪普遍采用\textbf{扫频超外差}结构：被测信号先经衰减器与预选滤波器进入第一混频器，与本地振荡器（LO）信号混频得到固定中频（IF）；LO 在处理器控制下按“扫频”方式改变频率，使输入频段内的各个频率分量依次被搬移到同一中频通带内，再经中频放大器、分辨率滤波器与检波器，最终在屏幕上绘出“幅度—频率”曲线。输入频率 $f_{\mathrm{in}}$、本振频率 $f_{\mathrm{LO}}$ 与中频 $f_{\mathrm{IF}}$ 满足
\[
    f_{\mathrm{IF}} = |\, f_{\mathrm{LO}} - f_{\mathrm{in}} \,| ,
\]
当 LO 扫过整个频段，整段频谱即被顺序测量。超外差结构借助中频滤波获得很高的灵敏度与动态范围，但也带来镜像响应与中频泄漏等问题，需由预选滤波与多次变频加以抑制。

\textbf{分辨率带宽（RBW）}与\textbf{视频带宽（VBW）}是理解频谱仪的两把钥匙。RBW 即中频通带的等效噪声带宽，它决定相邻谱线的可分辨程度：RBW 越窄，频率分辨力越高，但扫描同样频段所需时间越长。由于底噪功率正比于 RBW，将 RBW 由 $B_1$ 减小到 $B_2$ 可使显示平均噪声电平（DANL）降低
\[
    \Delta L = 10\log_{10}\!\left(\frac{B_1}{B_2}\right)\ \mathrm{dB} ,
\]
例如 RBW 由 $1\,\mathrm{MHz}$ 缩至 $1\,\mathrm{kHz}$，底噪下降约 $30\,\mathrm{dB}$，从而能发现更微弱的信号。VBW 则是检波后低通滤波器的带宽，用以平滑迹线、抑制显示抖动，通常取不大于 RBW。

\subsection{主要功能}

\begin{itemize}
    \item \textbf{频率分析}：分析信号的频率成分。可读取载波频率、偏移、占用带宽，定位异常谱线。
    \item \textbf{幅度测量}：测量各频率分量的幅度。以 $\mathrm{dBm}$ 或 $\mathrm{dB}\mu\mathrm{V}$ 等相对参考电平的单位表示绝对或相对功率。
    \item \textbf{谐波分析}：分析信号的谐波成分。对放大器、变频器做谐波失真评估，读取各次谐波相对基波的抑制比。
    \item \textbf{干扰检测}：检测信号中的干扰成分。在 EMC 预兼容测试中查找辐射与传导干扰源。
    \item \textbf{调制分析}：分析调制信号的特性。配合解调与专用选件，可观测带宽、邻道泄漏比（ACLR）等。
\end{itemize}

当被测信号极其微弱、逼近仪器底噪时，可开启\textbf{内置前置放大器}以提升灵敏度——前置放的增益 $G$ 与噪声系数 $F_{\mathrm{pre}}$ 会改善系统底噪，但前提是前置放自身的噪声系数要远低于后续链路。系统噪声系数（Friis 公式）为
\[
    F_{\mathrm{total}} = F_1 + \frac{F_2-1}{G_1} + \frac{F_3-1}{G_1 G_2} + \cdots ,
\]
可见第一级（前置放）的噪声系数与增益对整体灵敏度起决定性作用。需注意前置放也会放大输入端的强信号，可能使混频器过载，故应配合衰减器谨慎使用。

\subsection{应用领域}

\begin{itemize}
    \item \textbf{无线通信}：分析射频信号。测量发射机输出频谱、占用带宽、杂散与相位噪声。
    \item \textbf{音频测试}：分析音频信号的频率响应。配合低失真源评估放大器与扬声器的频响与失真。
    \item \textbf{电磁兼容}：检测电磁干扰。按标准限值扫描辐射与传导发射，定位超标频点。
    \item \textbf{雷达测试}：分析雷达信号。观测 chirp 线性度、脉冲频谱与杂散。
\end{itemize}

下表列出频谱仪使用中最关键的设置项及其影响。

\begin{table}[htbp]
    \centering
    \caption{频谱分析仪关键参数与影响}
    \begin{tabular}{|p{3.2cm}|p{4.4cm}|p{6.4cm}|}
        \hline
        \textbf{参数} & \textbf{符号/取值} & \textbf{对测量的影响} \\
        \hline
        \textbf{频率范围} & $f_{\mathrm{span}}$ & 屏幕覆盖的频段宽度 \\
        \hline
        \textbf{分辨率带宽} & RBW & 越小分辨力越高、底噪越低、扫描越慢 \\
        \hline
        \textbf{视频带宽} & VBW & 平滑迹线、抑制显示抖动 \\
        \hline
        \textbf{参考电平} & Ref Level & 决定纵坐标满刻度，须避免过载削顶 \\
        \hline
        \textbf{检波方式} & 峰值/均方根/采样 & 决定迹线对噪声与脉冲的表征 \\
        \hline
    \end{tabular}
\end{table}

\section{网络分析仪}
\label{sec:network_analyzer}

网络分析仪是用于测量网络参数的仪器。这里的“网络”并非计算机网络，而是指具有端口的线性微波/射频部件（如滤波器、放大器、天线、电缆）。网络分析仪通过向被测件（DUT）注入已知激励并测量其端口响应，定量描述 DUT 的反射与传输特性，是射频部件研发与生产的必备仪器。

\subsection{工作原理}

网络分析仪通过测量入射波和反射波来确定网络的S参数。采用“散射参数（S 参数）”描述端口关系最为方便：对于两端口网络，定义入射波 $a_1,a_2$ 与出射波 $b_1,b_2$，有
\[
    \begin{bmatrix} b_1 \\ b_2 \end{bmatrix}
    =
    \begin{bmatrix} S_{11} & S_{12} \\ S_{21} & S_{22} \end{bmatrix}
    \begin{bmatrix} a_1 \\ a_2 \end{bmatrix} ,
\]
其中 $S_{11}$ 为端口 1 的反射系数，$S_{21}$ 为正向传输系数（增益/插入损耗），$S_{12}$、$S_{22}$ 分别为反向隔离与端口 2 反射系数。矢量网络分析仪（VNA）同时测出各 S 参数的幅度与相位，故又称“矢量”分析。为实现精确测量，必须先用标准件进行\textbf{校准}，以消去测试电缆、连接器与仪器自身误差（方向性、源匹配、反射跟踪、传输跟踪等），常用方法为 SOLT（Short-Open-Load-Thru，短路—开路—负载—直通）。

\subsection{主要功能}

\begin{itemize}
    \item \textbf{S参数测量}：测量网络的散射参数。可同时给出幅度与相位，构成完整的频域模型。
    \item \textbf{阻抗测量}：测量网络的输入输出阻抗。由 $S_{11}$ 换算得到端口阻抗或导纳，判断与系统特性阻抗（通常 $50\,\Omega$）的匹配程度。
    \item \textbf{驻波比测量}：测量驻波比。由反射系数 $\Gamma$ 得电压驻波比
    \[
        \mathrm{VSWR} = \frac{1+|\Gamma|}{1-|\Gamma|} ,
    \]
    理想匹配时 $|\Gamma|=0$、VSWR = 1；反射越强，VSWR 越大。
    \item \textbf{增益测量}：测量网络的增益。由 $|S_{21}|$ 换算插入增益/损耗（dB），并观察随频率的变化即频响。
    \item \textbf{相位测量}：测量信号的相位变化。相位信息对滤波器群延时、放大器稳定性分析不可或缺。
\end{itemize}

与反射系数对应的\textbf{回波损耗（Return Loss, RL）}定义为
\[
    \mathrm{RL} = -20\log_{10}|\Gamma| \ \mathrm{dB} ,
\]
回波损耗越大表示匹配越好（反射越小）。例如 $|\Gamma|=0.1$ 对应 RL $\approx 20\,\mathrm{dB}$、VSWR $\approx 1.22$。下表给出几组典型对应关系，便于现场快速估算。

\begin{table}[htbp]
    \centering
    \caption{反射系数、回波损耗与驻波比对照}
    \begin{tabular}{|p{4.0cm}|p{4.4cm}|p{4.6cm}|}
        \hline
        \textbf{反射系数 $|\Gamma|$} & \textbf{回波损耗 RL (dB)} & \textbf{驻波比 VSWR} \\
        \hline
        $0.05$ & $26.0$ & $1.11$ \\
        \hline
        $0.10$ & $20.0$ & $1.22$ \\
        \hline
        $0.20$ & $14.0$ & $1.50$ \\
        \hline
        $0.33$ & $9.5$ & $2.00$ \\
        \hline
        $0.50$ & $6.0$ & $3.00$ \\
        \hline
    \end{tabular}
\end{table}

\subsection{应用领域}

\begin{itemize}
    \item \textbf{天线测试}：测量天线的阻抗和辐射特性。由 $S_{11}$ 评估驻波与带宽，配合其他手段评估效率。
    \item \textbf{滤波器测试}：测量滤波器的频率响应。读取通带插入损耗、带外抑制与群延时。
    \item \textbf{放大器测试}：测量放大器的增益和稳定性。在设定功率下测得 $S_{21}$ 与输入输出匹配，注意避免大信号烧毁端口。
    \item \textbf{传输线测试}：测量传输线的特性阻抗。定位电缆中的开路、短路与阻抗不连续点（时域反射 TDR）。
\end{itemize}

\section{逻辑分析仪}
\label{sec:logic_analyzer}

逻辑分析仪是用于分析数字信号的仪器。与示波器关注波形的模拟细节不同，逻辑分析仪只判定每一条线在采样时刻是“高”还是“低”（0 或 1），从而以多通道、长存储的方式重建数字系统的时序与协议，是调试微处理器总线、FPGA 与串行通信的利器。

\subsection{工作原理}

逻辑分析仪同时采集多个数字信号，并显示其时序关系。其内部由多路采样器构成：每个通道经阈值比较器把模拟电平判决为逻辑值，再按统一时钟（同步采样）或各通道跳变（异步采样）写入深存储器，最后以“时序图”形式在屏幕上回放。同步采样由外部或内部时钟驱动，适合分析总线协议；异步采样以远高于信号速率的时钟（通常数倍至数十倍）捕获，适合捕捉毛刺与不规则事件。采样必须满足
\[
    f_{\mathrm{sample}} \ge 2 f_{\mathrm{signal\ max}} ,
\]
实际为保证边沿不被漏采，常取信号最高速率的 $4\sim10$ 倍。

\subsection{主要功能}

\begin{itemize}
    \item \textbf{时序分析}：分析数字信号的时序关系。以多条波形对照观察建立/保持时间、延时与竞争冒险。
    \item \textbf{状态分析}：分析数字系统的状态变化。将总线按时钟采样为“状态序列”，按汇编或协议格式显示。
    \item \textbf{协议分析}：分析数字通信协议。内置 I2C、SPI、UART、CAN 等解码器，将比特流翻译为可读的帧与字段。
    \item \textbf{毛刺检测}：检测信号中的毛刺。利用硬件比较在单个采样周期内识别窄脉冲，避免被慢速采样所遗漏。
    \item \textbf{触发功能}：根据特定条件触发采集。可设置边沿、码型、序列、超时等多种触发，使分析仪在关键事件前后抓取数据。
\end{itemize}

\textbf{触发}是逻辑分析仪的灵魂：当被测系统运行极快、数据量巨大时，只有依靠精确的触发条件（如某个地址出现、某帧同步字到达），才能把感兴趣的一小段截取出来分析，否则海量采样将淹没有效信息。

\subsection{应用领域}

\begin{itemize}
    \item \textbf{数字电路调试}：调试数字逻辑电路。定位组合与时序逻辑中的错误与竞争。
    \item \textbf{微处理器测试}：测试微处理器的运行状态。观测地址/数据总线与外部存储、外设的交互。
    \item \textbf{通信协议测试}：测试串行通信协议。对 I2C、SPI、UART 等总线做物理层与协议层验证。
    \item \textbf{FPGA验证}：验证FPGA设计的正确性。对照设计预期检查内部逻辑端口的实际行为。
\end{itemize}

\section{功率分析仪}
\label{sec:power_analyzer}

功率分析仪是用于测量电功率的仪器。与仅测电平的功率计不同，功率分析仪面向交流系统的完整功率计量，能够同时处理电压与电流的波形，给出有功、无功、视在功率及功率因数，是电源、电机、变流器与节能评估中的关键仪表。

\subsection{主要功能}

\begin{itemize}
    \item \textbf{功率测量}：测量有功功率、无功功率、视在功率。对电压 $u(t)$ 与电流 $i(t)$ 同步采样后，有功功率
    \[
        P = \frac{1}{T}\int_0^T u(t)\,i(t)\,dt ,
    \]
    视在功率 $S = U_{\mathrm{rms}} I_{\mathrm{rms}}$，无功功率 $Q = \sqrt{S^2 - P^2}$（单相单频近似）。
    \item \textbf{能量测量}：测量电能消耗。对功率随时间积分得到 kWh 等能量量，支持长时间累计。
    \item \textbf{功率因数测量}：测量功率因数。定义为 $\cos\varphi = P/S$，反映有功功率在视在功率中所占比例。
    \item \textbf{谐波分析}：分析电流和电压的谐波成分。按 FFT 或离散傅里叶分解，给出各次谐波含量与总谐波畸变率（THD）。
    \item \textbf{效率测量}：测量设备的能源效率。同时测得输入与输出功率，计算 $\eta = P_{\mathrm{out}}/P_{\mathrm{in}}$。
\end{itemize}

对于含谐波的非正弦系统，功率分析仪依据 IEEE 1459 等方法处理，区分基波功率与谐波功率，避免简单相乘带来的误差。在高频开关电源与电机驱动中，电流波形含大量开关纹波，必须使用足够高的采样率与带宽，否则有功与谐波测量将失真。

\subsection{应用领域}

\begin{itemize}
    \item \textbf{电源测试}：测试电源的效率和稳定性。读取转换效率、纹波与动态响应下的功率变化。
    \item \textbf{电机测试}：测试电机的功率消耗。评估起动、额定与轻载工况下的输入功率与功率因数。
    \item \textbf{节能监测}：监测设备的能源消耗。对产线或楼宇做分项计量，定位能耗热点。
    \item \textbf{电力系统}：测量电力系统的功率参数。在电网与变流器中分析有功、无功潮流与谐波畸变。
\end{itemize}

功率分析仪常与\textbf{平均功率}与\textbf{峰值功率}概念并用：前者反映稳态能量，后者反映瞬态冲击。对于射频大功率测量，则多用基于定向耦合器的\textbf{功率计}，由耦合出的一小部分功率经二极管或热敏检波得到平均功率，并可在脉冲调制下借助包络检波测得峰值功率，二者互补，覆盖从工频到微波的广阔功率测量需求。

专业测试仪器的共同特点是高动态范围、宽频带与严格的校准，使用门槛也高于常用仪器。以频谱仪为例，读取微弱信号前应先设好参考电平与输入衰减，避免大信号使第一级混频过载而产生“假响应”；扫宽（Span）与分辨率带宽 RBW 的选择需在分辨力与扫描时间之间权衡，RBW 越窄底噪越低但扫描越慢。网络分析仪测量前务必完成校准（常用 SOLT），并定期复核以抵消电缆漂移；被测件与端口间应保持良好匹配，否则测得的反射将混入电缆误差。逻辑分析仪则要根据总线速率设置足够高的采样率，并合理配置触发，才能在海量数据中抓住关键事件。

仪器间的协同可构成完整的测试系统。例如评估一台射频发射机：用频率计或频谱仪测载频准确度，用功率计或频谱仪测输出功率，用矢量网络分析仪测天馈驻波，用频谱仪测输出频谱与杂散。只有各仪器量值均可追溯、设置互相自洽，测试结果才具有可比性。

选型时应重点核对：频率与带宽范围是否覆盖被测信号；动态范围与底噪是否满足最弱信号检测；幅度与相位的绝对准确度及稳定度；端口类型（N、SMA、BNC）与阻抗；所需通信接口（LAN、USB、GPIB）。切勿仅凭“最高指标”采购，而要看全系统误差预算中的瓶颈环节。专业仪器价格昂贵且对环境敏感，使用规范尤为重要：开机预热至稳定、定期自校与外接校准、保持端口清洁并以恰当扭矩连接（射频连接器不宜反复过度拧紧）、在静电与电磁受控场所操作。良好习惯能显著延长校准有效期并降低故障率。

以滤波器测试为例说明操作要点：矢量网络分析仪校准后连接滤波器，测 $S_{21}$ 得到插入损耗与通带纹波，测 $S_{11}$ 得到回波损耗或驻波比，在阻带读取带外抑制。若带外抑制不足，应先检查校准是否过期、电缆与连接器是否完好清洁，再怀疑被测件本身——这正是“先仪器后对象”的排错顺序。

数据表述方面，专业测试结果多以 $\mathrm{dBm}$、$\mathrm{dBc}$、$\mathrm{dB}$ 等对数单位给出。功率 $P$ 与 $\mathrm{dBm}$ 的换算为
\[
    P[\mathrm{dBm}] = 10\log_{10}\!\left(\frac{P}{1\,\mathrm{mW}}\right),
\]
相对量 $\mathrm{dBc}$ 表示相对于载波的电平。熟悉这些单位，有助于快速判断被测件是否达到设计指标。

在电磁兼容（EMC）测试中，频谱分析仪配合近场探头或天线可定位辐射热点，按标准限值判断合格与否；这类“预兼容”测试能在产品送检前发现大多数问题，降低整改成本。对于通信设备的研发与认证，专业测试仪器已从单台工具演变为覆盖“射频—基带—协议”的全套测量平台，熟练运用它们是通信设备工程师的核心能力之一。

专业仪器的\textbf{日常管理}可归纳为一份简易清单：每日检查自校/自检是否通过、风扇噪声与温度是否正常；每周清洁端口与通风口、确认校准有效期；每月做期间核查并比对历史数据；每年安排正式校准并更新台账。把清单制度化，比依赖个人记忆更可靠。

\textbf{安全}在专业仪器上更需重视：射频端口可能输出较大功率，误碰或误接会造成烫伤或器材损坏；高压电源与示波器高压差分探头须按规程使用；打开机壳前必须断电放电。任何维护不得破坏原有的安全设计与屏蔽完整性，否则可能引入新的干扰甚至触电隐患。

就\textbf{误差预算}而言，专业测量常以 $\mathrm{dB}$ 为单位分配：例如要求整机增益测量不确定度优于 $0.3\,\mathrm{dB}$，则网络分析仪的幅度不确定度、电缆与连接器的重复性、校准残余误差应各自留有余量并合成验证。当链路中任一项接近总限时，应优先升级该环节而非整体提高指标。

这类仪器的\textbf{校准周期}通常短于常用表：频谱仪、网络分析仪的建议周期多为 1 年，频率标准可半年；使用频繁或处于恶劣环境的应缩短。校准内容除幅度与频率准确度外，还包括噪声系数、相位线性度、端口匹配等专项，必要时由具备资质的机构实施。

掌握专业测试仪器，意味着既懂其背后的超外差、S 参数、采样与功率检波原理，也懂如何校准、如何解读不确定度、如何协同组成测试系统。唯有如此，才能把“看见信号”升级为“读懂信号”，为通信设备的研发、认证与运维提供坚实的数据支撑。
